% ============================================================
% CAPÍTULO 6 — Lecciones Aprendidas
% ============================================================
\chapter{Lecciones Aprendidas}
\label{cap:lecciones}

Este cap\'itulo recoge las lecciones aprendidas durante el ciclo de vida
del Proyecto Demeter, organizadas en tres dimensiones: t\'ecnicas,
organizativas y de gesti\'on. Estas lecciones constituyen un activo
adicional del proyecto, aplicables tanto a futuros proyectos de la
asociaci\'on como a otras iniciativas docentes similares.

\section{Lecciones t\'ecnicas}

\begin{itemize}[leftmargin=*]
    \item \textbf{Modelos compartidos como infraestructura central}:
    el uso de modelos Pydantic compartidos entre todas las capas para
    serializaci\'on/deserializaci\'on JSON y binaria result\'o
    especialmente acertado. La misma definici\'on de mensaje sirve
    para validaci\'on en backend y para empaquetado en firmware sin
    duplicaci\'on de l\'ogica.

    \item \textbf{Filosof\'ia de fallo silencioso en redes RF}: descartar
    silenciosamente las tramas corruptas evita tormentas de
    retroalimentaci\'on en presencia de interferencias. La l\'ogica de
    aplicaci\'on aporta los \texttt{ACK} necesarios.

    \item \textbf{Cuadros de \textit{Decisi\'on de Dise\~no} en
    documentaci\'on}: documentar expl\'icitamente las justificaciones de
    decisiones t\'ecnicas resulta cr\'itico para la mantenibilidad y la
    incorporaci\'on de nuevos miembros. Las cajas con razonamiento
    impl\'icito han demostrado ser una de las mejores aportaciones
    documentales.

    \item \textbf{Containerizaci\'on Docker en Raspberry Pi}: si bien
    introduce overhead inicial, la consistencia de despliegue y la
    facilidad de actualizaciones justific\'o ampliamente la decisi\'on.

    \item \textbf{Lecci\'on costosa: ausencia de cifrado en capa de
    campo}: planificar AES-128 como mejora futura en lugar de
    implementarlo desde el principio result\'o ser deuda t\'ecnica
    significativa al considerar despliegues en entornos de
    investigaci\'on con datos sensibles.

    \item \textbf{Patr\'on que se repetir\'ia: el modelo compartido
    entre capas}. Definir el protocolo una sola vez con modelos Pydantic
    en \texttt{Software/Common} y consumirlos desde el backend y desde la
    Raspberry Pi elimin\'o de ra\'iz toda una clase de errores de
    desincronizaci\'on entre extremos. Lo mismo aplica al patr\'on
    \textit{Strategy} sobre la interfaz \texttt{IComms} del firmware, que
    permiti\'o probar la l\'ogica de protocolo en el PC sin hardware.

    \item \textbf{Decisi\'on que se cambiar\'ia: medir antes de
    optimizar y antes de cerrar}. El sistema se dio por bueno sin una
    l\'inea base de latencia, p\'erdida de tramas ni autonom\'ia real de
    los nodos. Sin esas medidas no hay forma de demostrar que funciona ni
    de detectar una regresi\'on, y a\~nadirlas al final ya no es posible.
    La instrumentaci\'on m\'inima --- cobertura de tests y cuatro
    m\'etricas de campo --- deber\'ia haber existido desde la primera
    versi\'on.

    \item \textbf{Decisi\'on que se cambiar\'ia: validar el hardware
    cr\'itico antes de comprarlo}. El \textit{spike} de validaci\'on de
    bombeo se escribi\'o pero nunca se ejecut\'o, y para cuando se
    adquiri\'o la instalaci\'on hidr\'aulica (900\,\euro) el plan de
    pruebas ya se hab\'ia quedado obsoleto: estaba redactado para bombas
    de 5\,V y el material comprado era de 12\,V. El resultado es material
    sin desembalar. Un \textit{spike} que no se ejecuta no reduce ning\'un
    riesgo; solo documenta que se era consciente de \'el.

    \item \textbf{Decisi\'on que se cambiar\'ia: higiene del repositorio
    desde el primer d\'ia}. Versionar \texttt{node\_modules}, la
    documentaci\'on generada por Doxygen y ficheros de \textit{log}
    parece inocuo al principio, pero infla el repositorio en un orden de
    magnitud y destruye la posibilidad de medir nada autom\'aticamente
    sobre \'el. Un \texttt{.gitignore} correcto al inicio cuesta cinco
    minutos; corregirlo despu\'es implica reescribir el hist\'orico.
\end{itemize}

\section{Lecciones organizativas}

\begin{itemize}[leftmargin=*]
    \item \textbf{Distribuci\'on del conocimiento desde el inicio}: la
    concentraci\'on del desarrollo en un \'unico miembro fue la fuente
    principal de fragilidad del proyecto. Pr\'acticas como pair
    programming, code reviews sistem\'aticas y rotaci\'on de
    responsabilidades habr\'ian mitigado este \textit{single point of
    failure} estructural.

    \item \textbf{Documentaci\'on viva durante el desarrollo, no al
    final}: la documentaci\'on t\'ecnica generada al final del proyecto
    es valiosa, pero documentar progresivamente durante el desarrollo
    habr\'ia facilitado la incorporaci\'on de nuevos miembros y la
    resoluci\'on m\'as r\'apida de incidencias.

    \item \textbf{Onboarding sistematizado para nuevos miembros}: la
    asociaci\'on no contaba con un proceso definido de incorporaci\'on
    al proyecto. Una gu\'ia de \textit{onboarding} habr\'ia reducido la
    curva de aprendizaje de los miembros que se incorporaron en fases
    intermedias.

    \item \textbf{Distinci\'on entre documentaci\'on para evaluadores y
    documentaci\'on para mantenedores}: son audiencias y prop\'ositos
    distintos. La documentaci\'on para evaluadores enfatiza
    arquitectura y decisiones; la documentaci\'on para mantenedores
    enfatiza c\'omo hacer cosas concretas.
\end{itemize}

\section{Lecciones de gesti\'on}

\begin{itemize}[leftmargin=*]
    \item \textbf{Dimensionar las ambiciones al equipo realmente
    disponible}: el proyecto fue dise\~nado con una arquitectura de
    cuatro capas que requiere un equipo multidisciplinar consolidado.
    En retrospectiva, un MVP m\'as acotado, con menos capas pero
    iterables, habr\'ia sido m\'as compatible con la composici\'on real
    de la asociaci\'on.

    \item \textbf{MVP iterativo $>$ arquitectura aspiracional}: empezar
    con una versi\'on m\'inima funcional end-to-end y a\~nadir complejidad
    \'unicamente cuando se justifica con uso real es preferible a
    dise\~nar la arquitectura completa por anticipado.

    \item \textbf{Definir criterios de cierre desde el inicio}: cada
    proyecto deber\'ia tener criterios expl\'icitos de \'exito (cu\'ando
    se considera completado), pivote (cu\'ando cambiar de direcci\'on)
    y cierre (cu\'ando concluir formalmente). El Proyecto Demeter no
    los tuvo formalmente; la decisi\'on de cierre se toma ahora con
    criterios construidos a posteriori.

    \item \textbf{Reconocer cu\'ando cerrar es competencia profesional}:
    saber decir ``esto no podemos sostenerlo bajo las condiciones
    actuales'' es parte del aprendizaje en gesti\'on de proyectos. El
    presente informe es, en s\'i mismo, ejercicio de esa competencia.
\end{itemize}
